% !TEX root = ../main.tex
\pagestyle{fancy}
\addcontentsline{toc}{chapter}{Abstract}
\markboth{Abstract}{Abstract}
\newpage \thispagestyle{empty}

\section*{Abstract}
We develop a categorical framework for autonomous theorem proving by translating the problem of type inhabitation 
into a geometric graph exploration problem. Our main contributions are threefold.

First, we reformulate the functorial semantics of Generalized Algebraic Theories (GATs), 
leveraging Cartmell's Contextual Categories  to explicitly translate the syntactic problem of type inhabitation into 
a geometric graph search space. 
Second, we extend the Calculus of Constructions with a GAT signature and construct its term model.
Third, we design an Expansion Algorithm that exhaustively constructs the categorical search space of terms in both 
settings, and propose a Graph Neural Network-based Conjecturer-Prover architecture to guide and prune this search, overcoming 
the inherent combinatorial explosion, a severe limitation that we subsequently quantify by analyzing 
the super-exponential growth of the term search space in a foundational baseline of simply typed combinatory logic.

\section*{Resumen}
Desarrollamos las herramientas necesarias para abordar la demostraci\'on aut\'onoma de teoremas desde la teor\'ia de categor\'as,
traduciendo el problema de habitaci\'on de tipos en un problema de exploraci\'on geom\'etrica de grafos.
Nuestras contribuciones principales son las siguientes:

En primer lugar, reformulamos la sem\'antica functorial de las Teor\'as Algebraicas Generalizadas (GATs),
utilizando las Categor\'ias Contextuales de Cartmell para traducir expl\'icitamente el problema sint\'actico de habiitaci\'on de tipos en un 
problema geom\'etrico de exploraci\'on de grafos.
En segundo lugar, extendemos el C\'alculo de Construcciones con una signatura GAT y construimos su modelo sintactico.
En tercer lugar, dise\~namos un Algoritmo de Expansi\'on que construye 
de forma exhaustiva el espacio de b\'usqueda categ\'orico de t\'erminos en ambos entornos, y proponemos una 
arquitectura Conjeturador-Demostrador basada en Redes Neuronales de Grafos para guiar dicha busqueda,
superando la explosi\'on combinatoria inherente al problema. Cuantificamos este crecimiento superexponencial del espacio de b\'usqueda 
de t\'erminos en el entorno de la l\'gica combinatoria con tipos simples.

{\let\thefootnote\relax\footnote{2020 Mathematics Subject Classification. 03B38, 18C10, 68T30, 03F50}}